\subsection{有限基点型4・31の全係数被覆とH3排除}
\label{subsec:cubic-CY-finite43-cover}
\begin{theorem}\label{thm:cubic-CY-finite43-empty}
零交換子CY包絡を持つH3の対称射影の基点型は、$4,31$ のどちらにもならない。
従ってこの領域の必要対称射影型は $1111,211,22$ に限られる。
\end{theorem}
\begin{proof}
固定出典\cite{AS4CYpinned}の曲線的基点型4・31の全係数定理から、必要被覆だけを用いる。
両方の全21混合方向、印付き射影の全安定化群、非零係数による尺度変更と
定数非零pivotによる平行移動を保持する。
上流の全検証器を固定referenceから隔離した場所で再実行し、以下に用いる全必要分岐を確認した。
当repoではさらに元tensorと全空ideal・全非正則族・残る三表示への関係移送を独立に再生した。
AS4の十分性、Koszul filtrationや全自己同型の結論はこの排除に用いない。

\paragraph{一点長さ4の型}
$U=\langle ac-b^2,ad-bc,bd,c^2,cd,d^2\rangle$ とする。
基点環 $k[t]/(t^4)$ の全線型安定化群は、$t$ の局所自己同型と可逆元の乗算で尽くされる。
rank6の印付けを正規化した後も、固定出典の全四変数ユニポテント部分とトーラスを残す。
Schur条件を全21係数で扱うと、$x_0\ne0$ の二枝、次の $x_1\ne0$ の二枝は全て単位idealである。
ここでの各gaugeの主係数は非零定数であり、他の自由係数の特殊値を除いていない。
残る $x_0=x_1=x_4=x_5=x_9=0$ の下で、
$k=x_{10}-2x_6=0$ は実平方 $d^2$ を与えるので、$k=1$ へ正規化できる。

先頭四係数を $(u,v,w,s)=(x_2,x_6,x_7,x_{15})$ とすると、全必要式は
\[
\begin{gathered}
 v(-2s+u-v+2)=0,\quad s=(2v+1)(u+2),\quad w(u+w)=0,\\
 s(1-2w)-uv+uw-2vw+2w=0
\end{gathered}
\]
であり、幾何学的な全解は
\[
\begin{array}{c|rrrr|c}
 &u&v&w&s&\text{残余}\\\hline
0&-5/2&-1/2&5/2&0&\dim C_4=37\\
1&-2&0&0&0&\text{空}\\
2&-2&0&2&0&\text{空}\\
3&-5/3&-1/7&0&5/21&\mathcal C_4^{(1)}\\
4&-1&0&1&1&\mathcal C_4^{(2)}
\end{array}
\]
となる。$v=0$ と $v\ne0$ の場合分けで全てを得られ、
途中の分母 $4v+1=0$ は元式に反する。
なお、必要先頭idealが与える消去式は $v^2(2v+1)(7v+1)=0$ であり、
この五tupleの被覆を、必要schemeの被約性や相対分類と同一視してはいない。
残余の二空idealと三つの全orbit正規化を確認し、場合0を元六関係の四次商37次元で排除する。
場合3,4は定理\ref{thm:as4-curvilinear4-h3-exclusion}の二固定表示へ文字どおり移り、
中心三次が零なのでH3でない。

\paragraph{長さ3と1の型}
$U=\langle ac-b^2,bc,c^2,ad,bd,cd\rangle$ とする。
基点環 $k[t]/(t^3)\times k$ の全安定化群とrank6印付けを固定した全21係数から始める。
$x_0\ne0$ 又は $x_{10}\ne0$ の二枝は、全定数pivot消去後の単位idealで消える。
$x_9\ne0$ の全normal formを含め、その他の非正則な代表・族は次表で全て排除される。
\[
\begin{array}{c|c|c}
\text{枝}&\text{代表又は自由係数}&\text{商次元の下界}\\\hline
x_9\ne0&\text{有理一点と符号付き二直線}&\dim C_4\ge37\\
x_9\ne0&\mathbb Q(i,\sqrt2)\text{ の四点}&\dim C_4=36\\
x_{12}\ne0&h\ne0&\dim C_4\ge37\\
x_{12}\ne0&\text{符号付き二直線}&\dim C_5\ge62\\
x_{17}\ne0,\ x_1=0&\mu\ne0,-1,\ w&\dim C_4\ge39\\
x_{17}\ne0,\ x_1=0&\mu=-1,\ w,y&\dim C_4\ge39\\
x_{17}\ne0,\ x_1=0&\mu=0,\ r\ne0,t&\dim C_4\ge42
\end{array}
\]
ここには、$\mu=-1$ で平行移動gaugeが退化する例外を自由係数 $y$ を残して含めた。
この例外を一般 $\mu$ のchartへ無条件に特殊化してはいない。
二直線の五次下界62は、必要な $\binom83=56$ に反する。

各自由係数族では、元六関係の分母は明記した非零parameterの冪だけである。
一般分数体上の元全語関係行列のrankは、許容閉値への特殊化で増えない。
従って一般商次元の下界は全閉値に及ぶ。
元の全関係挿入を用いる独立な商構成で、表の全13代表・族を再計算した。
これは一般商基底を特殊化して平坦性を推測する議論ではない。

残余枝の被覆も省略しない。
$x_9=0$ では二つの三乗恒等式から $x_{11}=x_{15}=0$。
$x_{12}=0$ も残った場合、$x_{17}=0$ は実平方 $c^2$ を与える。
$x_{17}=1$ に正規化した全六残余式のうち、$x_1\ne0$ は一つの固定orbit
\[
 x_1=x_3=x_{14}=x_{17}=1,\quad x_2=x_4=-1/16,\quad g=2
\]
へ移り、これは定理\ref{thm:as4-curvilinear4-h3-exclusion}の $\mathcal C_{31}$ に一致する。
$x_1=0$ の全場合は表の最後の三行である。
従って全候補が消え、型31もH3を与えない。
\end{proof}

独立再現器は\\
\texttt{verify\_cubic\_CY\_finite43\_cover.py}。
固定証明の必要データだけを最小入力に保持し、二つの元21係数frame、全空ideal、
一点型の先頭消去式、全14非正則代表・族の元Hilbert、
三つの既存排除表示への定数可逆な関係基底移送を確認する。
非零gaugeの全被覆は固定出典の通常証明と上流全control再実行に依存し、
独立商計算・形式化・AS十分性を混同しない。
残る三射影型の全係数被覆と、他の交換子次元のCY中央切断は未完成である。
